% Slides — Capítulo 14: Fundamentos de Tempestades
\documentclass[aspectratio=169,11pt]{beamer}
\usepackage{../comum/temameteo}
\graphicspath{{../figuras/cap14/}}

\title[Cap. 14 --- Tempestades]{Capítulo 14 --- Fundamentos de
Tempestades}
\subtitle{Meteorologia Prática}
\author{}
\date{}

\begin{document}

\begin{frame}
  \titlepage
  \vfill
  \atribuicao
\end{frame}

\begin{frame}{Roteiro}
  \tableofcontents
\end{frame}

% ---------------------------------------------------------------
\section{A receita de quatro ingredientes}

\begin{frame}{O objeto do capítulo}
  \begin{columns}
    \column{0.4\textwidth}
    \begin{itemize}
      \item Cumulonimbus: a máquina térmica de 30--60\,min que
            concentra a energia de um dia inteiro de sol.
      \item Receita: \alert{umidade $+$ instabilidade $+$ gatilho $+$
            cisalhamento}.
      \item Os três primeiros decidem SE; o quarto decide QUAL.
    \end{itemize}
    \column{0.6\textwidth}
    \figlivroref{Fig.~14.1}{cumulonimbus com bigorna: o produto final}{Gene Rhoden / weatherpix.com}
  \end{columns}
\end{frame}

% ---------------------------------------------------------------
\section{CAPE e CIN}

\begin{frame}{A energia armada no Skew-$T$}
  \begin{columns}
    \column{0.4\textwidth}
    \eqdestaque{$\mathrm{CAPE} = \displaystyle\int_{LFC}^{EL}
      g\frac{T_{v,p}-T_{v,e}}{T_{v,e}}dz$}
    \vspace{0.4em}
    \begin{itemize}
      \item Área positiva $=$ CAPE; negativa (a tampa) $=$ CIN.
      \item LCL $\to$ LFC $\to$ EL: as três alturas da parcela
            (Caps.~4--5).
      \item $<$500 fraca $\cdot$ 500--1500 moderada $\cdot$
            $>$2500\,J/kg severa.
    \end{itemize}
    \column{0.6\textwidth}
    \figlivroslide{fig_14_25.jpg}{0.56\textheight}{Fig.~14.25 ---
      sondagem pré-tempestade: CAPE, CIN e a tampa}
  \end{columns}
\end{frame}

\begin{frame}{Do CAPE ao updraft — e o papel do CIN}
  \begin{columns}
    \column{0.5\textwidth}
    \eqdestaque{$w_{max} = \sqrt{2\,\mathrm{CAPE}}$\quad
      (real $\approx$ metade)}
    \vspace{0.4em}
    \begin{itemize}
      \item CAPE 2000 $\Rightarrow$ 63\,m/s teórico, $\sim$30 real:
            entranhamento $+$ peso da água cobram (T1, N2).
      \item CIN moderado $=$ panela de pressão: acumula o dia e
            \alert{explode à tarde} (N1).
      \item Gatilhos: frente, brisa, linha seca, montanha, rajada da
            vizinha (Caps.~12, 17).
    \end{itemize}
    \column{0.5\textwidth}
    \figlivroslide{fig_14_27a.jpg}{0.52\textheight}{Fig.~14.27 --- a
      tarde esquenta, a tampa derrete, a convecção dispara}
  \end{columns}
\end{frame}

% ---------------------------------------------------------------
\section{O ciclo de vida}

\begin{frame}{Três estágios, meia hora cada}
  \begin{columns}
    \column{0.25\textwidth}
    \begin{itemize}
      \item Torre (só updraft);
      \item maduro (up $+$ down, chuva forte);
      \item dissipante (só downdraft).
      \item A chuva que cai DENTRO do updraft o mata: célula única
            se suicida.
    \end{itemize}
    \column{0.375\textwidth}
    \figlivroslide{fig_14_6b.jpg}{0.4\textheight}{Fig.~14.6 --- torre
      de cumulus e estágio maduro}
    \column{0.375\textwidth}
    \figlivroslide{fig_14_6d.jpg}{0.4\textheight}{Fig.~14.6 ---
      dissipação: o downdraft vence}
  \end{columns}
\end{frame}

\begin{frame}{A piscina fria e a frente de rajada}
  \begin{columns}
    \column{0.4\textwidth}
    \eqdestaque{$c \simeq \sqrt{2gH\,\Delta\theta_v/\theta_v}
      \sim 15$--$25$\,m/s}
    \vspace{0.4em}
    \begin{itemize}
      \item Chuva evapora no ar seco sub-nuvem $\to$ ar frio desaba e
            se espalha (corrente de densidade, T4).
      \item A frente de rajada ergue o vizinho: célula-filha —
            multicélulas se auto-propagam.
      \item Linhas duradouras: RKW $c/\Delta u \approx 1$ (N4).
    \end{itemize}
    \column{0.6\textwidth}
    \figlivroslide{fig_14_12.jpg}{0.5\textheight}{Fig.~14.12 --- a
      frente de rajada avançando sob a tempestade}
  \end{columns}
\end{frame}

\begin{frame}{Multicélulas e linhas de instabilidade}
  \begin{columns}
    \column{0.4\textwidth}
    \begin{itemize}
      \item Trem de células em estágios diferentes; a rajada comum
            dispara cada filha.
      \item Linha de instabilidade: a versão organizada, com jato de
            retaguarda e \emph{bookend vortices}.
      \item No radar: eco em arco (bow echo) $=$ rajadas destrutivas.
    \end{itemize}
    \column{0.6\textwidth}
    \figlivroslide{fig_14_14.jpg}{0.5\textheight}{Fig.~14.14 --- corte
      de linha de instabilidade: piscina fria e jato de retaguarda}
  \end{columns}
\end{frame}

% ---------------------------------------------------------------
\section{Cisalhamento e supercélulas}

\begin{frame}{O hodógrafo decide}
  \begin{columns}
    \column{0.4\textwidth}
    \begin{itemize}
      \item $<$10\,m/s (0--6\,km): célula única; 10--20: multicélula;
            $>$20 $+$ curvatura: \alert{supercélula}.
      \item Cisalhamento inclina o updraft: a chuva cai FORA dele —
            a tempestade não se afoga.
      \item Curvatura do hodógrafo $\to$ helicidade $\to$ updraft
            \alert{rotativo} (mesociclone) — SRH no N3.
    \end{itemize}
    \column{0.6\textwidth}
    \figlivroslide{fig_14_64b.jpg}{0.52\textheight}{Fig.~14.64 ---
      hodógrafo e a helicidade relativa à tempestade (área
      sombreada)}
  \end{columns}
\end{frame}

\begin{frame}{A supercélula por dentro}
  \begin{columns}
    \column{0.4\textwidth}
    \begin{itemize}
      \item Updraft rotativo inclinado (20--40\,m/s), bigorna com
            \emph{overshooting top} (Cap.~5).
      \item Downdrafts FFD e RFD; wall cloud sob o mesociclone
            (Cap.~15: tornadogênese).
      \item Horas de vida: a estrutura em equilíbrio dinâmico.
    \end{itemize}
    \column{0.6\textwidth}
    \figlivroslide{fig_14_4a.jpg}{0.56\textheight}{Fig.~14.4 ---
      anatomia da supercélula}
  \end{columns}
\end{frame}

\begin{frame}{No radar: o eco em gancho}
  \begin{columns}
    \column{0.4\textwidth}
    \begin{itemize}
      \item O mesociclone enrola a chuva: gancho na refletividade
            (Cap.~8) $+$ dipolo de velocidade Doppler.
      \item Região de eco fraco (WER): o updraft tão forte que a
            chuva não cai dentro.
      \item BRN $= \mathrm{CAPE}/(\tfrac12\Delta U^2)$ resume:
            supercélulas em 10--50 (T5).
    \end{itemize}
    \column{0.6\textwidth}
    \figlivroslide{fig_14_18a.jpg}{0.52\textheight}{Fig.~14.18 ---
      assinatura clássica de supercélula no radar}
  \end{columns}
\end{frame}

% ---------------------------------------------------------------
\section{Síntese e laboratório}

\begin{frame}{O mapa do capítulo}
  \eqdestaque{umidade $+$ CAPE $+$ gatilho $+$ cisalhamento;\quad
    $w = \sqrt{2CAPE}$;\quad $c = \sqrt{2gH\Delta\theta_v/\theta_v}$;
    \quad BRN}
  \vspace{0.6em}
  A energia é térmica (Caps.~3--5), a organização é do vento
  (Cap.~11): célula única $\to$ multicélula $\to$ supercélula conforme
  o hodógrafo. Os produtos perigosos: Cap.~15.
\end{frame}

\begin{frame}{Laboratório numérico --- notebook do Cap.~14}
  \texttt{notebooks/cap14\_tempestades.ipynb}
  \begin{enumerate}
    \item CAPE/CIN com MetPy (Skew-$T$ sombreado).
    \item $w_{max}(\mathrm{CAPE})$ e o fator 1/2 real.
    \item Hodógrafos: única × multi × supercélula.
    \item Piscina fria e a EDO do downdraft.
  \end{enumerate}
  \vspace{0.4em}
  \textbf{Exercícios}: T1--T5 e N1--N4 nas notas; células reservadas no
  notebook. Sugestão: rodar o Caso~1 com a sondagem do dia e
  ``prever'' a tarde com a turma.
  \vfill
  \atribuicao
\end{frame}

\end{document}
